The \texttt{with-in} statement defines \emph{bounded variable}s, which can then be iterated over using the \texttt{foreach} statement. The basic syntax is as follows:

\begin{lstlisting}[language=croqtile,caption={Indexing one shape}]
with index in [6] {
  foreach index {
    // Perform operations with each index
  }
}
\end{lstlisting}

In this example, the \texttt{foreach} block iterates 6 times, with the value of the \textbf{iteration variable} \texttt{index} ranging from \texttt{0} to \texttt{5} incrementally.

It is also possible to iterate over \emph{bounded ituple}s. For example:

\begin{lstlisting}[language=croqtile,caption={A tuple of indices}]
with index = {x, y} in [6, 17] {
  foreach index { }
}
\end{lstlisting}

This code defines a \emph{bounded ituple} \texttt{index}. Iterating over \texttt{index} is equivalent to a nested loop structure:

\begin{lstlisting}[language={},caption={The generated C++ loop}]
for (int x = 0; x < 6; ++x)
  for (int y = 0; y < 17; ++y) { }
\end{lstlisting}

The left-to-right order of indices in a bounded tuple determines the outer-to-inner nesting of the loops. Keep this ordering in mind when the order of iteration matters.

The same outer-to-inner ordering applies to a comma-separated list of bounded variables:

\begin{lstlisting}[language=croqtile,caption={Several bounded groups}]
with index = {x, y} in [6, 17], iv in [128] {
  foreach iv, index { }
}
\end{lstlisting}

In this code, two bounded variables are defined. The following \texttt{foreach} block is equivalent to a multi-level nested loop:

\begin{lstlisting}[language={},caption={The generated inner loop}]
for (int iv = 0; iv < 128; ++iv)
  for (int x = 0; x < 6; ++x)
    for (int y = 0; y < 17; ++y) { }
\end{lstlisting}
